% ===================== PREFACE =====================
\cleardoublepage
\chapter*{Preface}
\addcontentsline{toc}{chapter}{Preface}
\markboth{Preface}{Preface}

\begin{center}
\begin{minipage}{0.88\textwidth}
\small\itshape\color{steel}
\centering
Begun 13 April 2018.
\end{minipage}
\end{center}
\vspace{10pt}

I started writing this book on the night of April 13th, 2018, because that was
the night Art Bell died.

The timing is not a coincidence and I am not going to dress it up as one. His
death is the reason this book exists.

I was in a computer lab at the University of Lethbridge when I found out. It
was late, the building was effectively empty, and I had a pair of headphones on
because I was doing what I had done in that lab a hundred nights before, which
was listening to \emph{Coast to Coast AM} while I pushed pixels around a
screen. What I was working on that night is the picture on the front of this
book --- the artwork is mine, made in that room, on that machine, and I am
telling you so here rather than burying it in a credits line, because of what
happened while I was making it. George Noory came out of a break and said, in
that unhurried way he has when the news is bad, that Art Bell was dead.

I sat in a room full of empty chairs and cried at a monitor. I want to be
precise about that, because it is the only genuinely honest sentence I can
offer about my own credentials. Not a colleague. Not a friend. A voice --- a
voice I had never once spoken to --- announced by the one voice in that studio
I had spoken to, because Noory had put me on the air the previous autumn to ask
what a Canadian university thought it was doing teaching this subject at all.
It landed like a death in the family, because that is
what thirty years of somebody keeping you company at three in the morning
actually builds. I thought about the fact that the voice which had kept me
company through a decade of insomnia had stopped, and that nobody was going to
pick up the phone at that hour and let a truck driver in Nevada describe what
he had just seen without laughing at him. By the time I got up from that chair
I had decided that somebody had better write the rest of it down. I opened a
blank document before I left the building, and the file I had been drawing when
Noory said it became the cover of the thing I opened it for.

It landed harder than it should have, and I know exactly why. Bell died one
month --- almost to the day --- after Stephen Hawking. I had spent March
reading obituaries of the man who taught me that you can reason your way to a
black hole without ever visiting one, and April listening to the man who taught
me that you let the witness finish his sentence. Losing both inside four weeks
felt less like two deaths than like the collapse of a single method: rigour at
one end, patience at the other, and nothing much left in the middle.

There is a detail I have never quite been able to leave alone. Hawking died on
a Wednesday. Wednesday evening was already the slot I had booked for the
lectures --- chosen for dull reasons of room availability
and nothing else --- and so every week I stood up to do this in front of
whoever had turned up on the same weeknight I had lost him. The night the
lectures reached the question of whether we could ever talk to something not
from here, I dedicated the hour to him out loud, before the first slide. That
dedication is still in this book. It is the opening of Chapter
\ref{ch:contact}, and it is the only part of the lecture series I have
transcribed word for word.

That was the thing about Bell. He let people talk. He did not endorse them, and
he did not humiliate them, and if you listen to enough of those tapes you
realise how genuinely rare that combination is. Somewhere between the
evangelists and the debunkers there was a man with a microphone in the high
desert who had figured out that the honest position is usually ``well, tell me
more.'' That is what I started typing towards that night, and somewhere in the
first hour of it, it occurred to me quite suddenly that the thing I actually
wanted to write about was not flying saucers at all.

\hi{This is not a UFO book.} I understand that it looks like one. It has the
word \emph{Ufology} on the cover in large letters, it has chapters on Roswell
and Rendlesham and Area 51, and if you shelve it anywhere it will be shelved
next to books with considerably worse research and considerably better sales.
But the subject of this book is not unidentified flying objects.
\hi{The subject of this book is our spectacular, ongoing, almost artistic
\mbox{\textbf{FAILURE TO COMMUNICATE!}}}

Consider what actually happens in a UFO case. A person sees something. That
something is, by definition, a thing they cannot identify --- so already the
signal has arrived damaged. They then have to convert it into words, and the
words available to them are terrible: ``saucer,'' ``craft,'' ``triangle,''
``it moved like nothing I've ever seen.'' Those words go to an investigator,
who has his own vocabulary and his own hopes. The investigator's report goes to
an agency, which has an interest in the answer. The agency's summary goes to a
journalist, who has a deadline. The journalist's article goes to the public,
which has already decided. By the time the story reaches you it has been
translated six times by six parties with six different agendas, and we then act
surprised that it does not survive scrutiny.

At no point in that chain does anyone lie, necessarily. That is what I find so
interesting. \hi{You do not need liars to produce a lie; you only need enough
people sincerely talking past one another.} Wrap that around a subject with a
stigma attached, so that half the witnesses will not speak and half the
scientists will not listen, and you have built a machine that reliably converts
real events into unusable noise. Seventy-five years of that machine running at
full capacity is roughly what we mean when we say ``the UFO phenomenon.''

So the questions that organise this book are communication questions. How is
the question worded, and what does the wording smuggle in? Who is required to
prove what, and who quietly got excused? What did the witness actually say,
before it was tidied up? What would count as evidence here, and has anyone
bothered to say so out loud? These are unglamorous questions. They are also the
only ones that have ever moved a case forward.

And there is a larger version of the problem waiting at the end of the book,
which is the part that genuinely keeps me up --- Bell's hours, appropriately.
If something out there ever does try to talk to us, the encounter will not fail
because of physics. \hi{It will fail because we are bad at listening, bad at
saying what we mean, and worse at admitting when we have understood nothing at
all.} We cannot reliably transmit the contents of one human night sky across a
single dinner table without it turning into an argument. That is the skill we
are short on. That is what I think is worth practising, and a disreputable
subject with a seventy-five-year paper trail turns out to be an excellent
place to practise it.

\hi{A word about how this book found out it had to be a book.}

The lectures came first. I started the U of L Ufology Group in the autumn of
2017, months before that night in the lab --- which is why Noory knew who I was
when the news came, and why the book has a lecture series behind it rather than
the other way round. It
was not an informal thing --- it was an officially registered club at the
University of Lethbridge, chartered through the Students' Union with all the
paperwork that implies: a constitution, an executive, a room booking, a table at
clubs week. I wanted a place on a real campus where this material could be
treated as a subject rather than as a punchline.

The response was, at first, ridiculous. Over a hundred people put their names on
the mailing list in the first week. I had not expected anything like it, and I
spent that week rewriting my plans upward, worrying about whether the room I had
booked was going to be big enough.

Six people came to the first lecture. Half a dozen. In a room I had worried
about filling.

I have thought about that gap --- a hundred names, six chairs --- more than is
probably healthy, and I no longer think it was a failure of interest. A hundred
people genuinely wanted this to exist. Signing a mailing list costs nothing and
says something true: yes, I am curious, yes, somebody should be doing this.
Walking into a lecture hall on a weeknight to be seen studying flying saucers
costs considerably more. The stigma described a few pages ago is not an abstract
sociological force. It is the exact distance between a hundred signatures and
six attendees, and I have measured it personally.

What settled in after that was two regulars --- two people who came to
everything, whose names belong in this preface more than mine does --- plus a
rotating cast who appeared for whichever topic had caught their eye. Somebody
would turn up for Roswell and never return. Somebody else would materialise
weeks later because they had heard I was doing time travel that night, sit
through it, and leave looking politely baffled.

And that is the thing that turned a lecture series into this textbook, because
those baffled faces were not the audience's fault. They were mine. If you came
only to the time-travel talk, you had missed the material that lecture was
standing on --- how a claim gets worded, who carries the burden of proof, what
an interesting answer looks like as opposed to a correct one, what the
documentary record actually consists of. Without that, ``time travel'' is just
an hour of entertaining nonsense delivered by a man with slides. With it, the
same hour is an exercise in bounded reasoning. Same words, same room, and
entirely different content depending on whether the listener had the previous
weeks in hand. I was watching, in real time and at close range, the exact
communication failure this whole book is about --- and I was the one causing it.

So: \hi{this book is built to be read in order, front to back, and that is not a
preference.} Every chapter assumes the chapters before it. The vocabulary is
built up deliberately, terms are defined once and then used, and the reasoning
tools installed early are applied without re-explanation later. If you open this
volume in the middle because a chapter title caught your eye, you will get the
same experience my drop-in students got: it will seem to make sense, you will
follow the sentences, and you will draw conclusions the material does not
support, because the part that constrained those conclusions was in a chapter
you skipped. A textbook cannot repeat itself twenty times to accommodate
arriving late. It can only tell you, once and plainly, that the order is load
bearing.

I want to be straight with you about where that structure comes from, because I
did not invent it and I did not borrow it from a textbook. I borrowed it from
the other side. Pick up the great mythology-builders of this field --- Corso on
Roswell, Cooper on the government, Icke on everything --- and you will notice
that none of them argues case by case. They build. Chapter one establishes a
premise, chapter four leans on it as though it were settled, chapter nine treats
the whole accumulated edifice as a foundation, and by the end the reader is
nodding along to a claim he would have laughed at three hundred pages earlier.
It is a genuinely powerful piece of engineering, and it is the reason those books
convert people while the debunkings mostly do not. The debunkings are essays. A
Shermer piece, a Nickell investigation, a Klass takedown --- excellent work, and
each one a self-contained hour that leaves the reader exactly where it found
him.

\actualq[problem]{the structure of a cumulative argument.}{what that structure
has been pointed at.}

A cumulative argument is simply what you build when each conclusion
legitimately depends on the one before it, and that describes rigour at least as
well as it describes mythology. So this book uses the believers' architecture and
aims it in the opposite direction: the thing that accumulates, chapter after
chapter, is not a hidden history of contact but a method --- how wording
predetermines answers, who owes proof to whom, what the documentary record
actually consists of, why an interesting explanation and a correct one are
different animals. By Chapter~\ref{ch:contact} you should be able to take a case
apart yourself, and the reason you will be able to is that every chapter before it
built the tool in your hands one piece at a time.

Which is also, and I want this understood plainly, why reading it out of order
does not merely cost you some background. \hi{It reverses the effect.} Every
cumulative structure carries the same risk, mine included: a chapter read on its
own feels persuasive, because the sentences are still doing their work and the
constraints that were supposed to bound them are sitting unread forty pages back.
That is exactly the mechanism by which the mythology-builders make a bad case
feel inevitable. Read this book in sequence and the accumulation makes you
harder to fool. Read it in fragments and it becomes one more volume that made
somebody confident about something he never checked --- and I would rather not
have written that book by accident.

\hi{Everything in this book was a lecture first.} The course material was written in 2017, in my
third year at the University of Lethbridge, for the club I had just founded, and it was written
under the only condition that matters for teaching: I had to stand in a room the following week
and say it out loud to people who were free to walk out. That is a harsher edit than any
manuscript review. A sentence that does not survive being spoken aloud to a dozen people in a
booked classroom does not belong in a textbook either.

At the end of that first year, in 2018, the University of Lethbridge Students' Union named the
U~of~L Ufology Group \emph{Outstanding New Club} --- the award given to the best new club on
campus, presented to the newly ratified club judged to have mounted the most outstanding and
creative programme of events over the year. I am aware of how that sounds in a book about flying
saucers, and I am going to say it anyway, because it is a matter of record and I have the plaque:
\hi{the course material this textbook is built on is award-winning.} A students' union is not a
scientific body and the award is not a peer review. What it is, is documentary evidence that a
campus with no stake whatsoever in ufology watched this material being taught for a year and
concluded it was the best new thing on it.

I ran the series a second time during the 2024--2025 school year, having revived the club after
enrolling at Lethbridge Polytechnic for a certificate in Virtual Reality studies. Between the two
runs I revised and updated the whole of the original course material --- seven years of new
documents, new hearings and new admissions had landed in the interval, including the ones
Section~\ref{sec:121617} treats as the dividing line of the entire subject, and a course written
before December 2017 could not honestly be re-delivered as though nothing had happened.

And now the material has been revised and updated a third time, with the help of AI, before being
made into this book. I would rather state that plainly on the way in than have it noticed on the
way out. What the machine was used for was labour and reach: checking dates, chasing primary
documents, locating public-domain illustration, holding a five-hundred-page structure in view at
once, and arguing with me about wording until the wording said what I actually meant. What it was
not used for is the argument. The method in this book, the judgements about which cases survive
and which do not, and every conclusion you are about to read are mine, and were mine when they
were lecture notes for a handful of students in Alberta.

% The award pair and the Open House photograph are each wrapped in one outer
% minipage, which cannot be split across a page break, so the pictures and the
% credit line that explains them always travel together. Without this the
% digital edition broke the pair away from its own caption.
\par\vspace{12pt}
\begin{center}
\begin{minipage}{\textwidth}
\centering
\begin{minipage}[c]{0.50\textwidth}
\centering
\includegraphics[width=\textwidth]{pf_award_cert.jpg}
\end{minipage}\hfill
\begin{minipage}[c]{0.44\textwidth}
\centering\includegraphics[width=\textwidth]{pf_award_photo.jpg}
\end{minipage}
\par\vspace{5pt}
\begin{minipage}{0.86\textwidth}
{\RaggedRight\centering\scriptsize\color{steel}\itshape Left: the citation itself --- ``Outstanding
New Club\ldots\ is hereby granted to The U of L UFOLogy Group,'' awarded for the 2017/2018 academic
year and signed by Students' Union president Hailey Babb. Right: collecting it. The award is given
to the newly ratified club judged to have run the most outstanding and creative programme of the
year, and the programme in question was the first draft of this book. Author's own
photographs.\par}
\end{minipage}
\end{minipage}
\end{center}
\vspace{10pt}

\begin{center}
\begin{minipage}{\textwidth}
\centering

\includegraphics[width=0.72\textwidth]{pf_polytech_table.jpg}\par
\vspace{5pt}
\begin{minipage}{0.78\textwidth}
{\RaggedRight\centering\scriptsize\color{steel}\itshape Seven years on and one institution across
town: the revived club's table at the Lethbridge Polytechnic Open House, autumn 2024, at the start
of the second run of the series. Back issues of the \emph{MUFON UFO Journal}, a copy of \emph{The
Roswell Incident}, the Lethbridge astronomy handouts, the Area~51 cap and the same row of little
grey figures that had sat on the table in 2017. The folding table does not change much. The
material had been rewritten from end to end. Author's own photograph.\par}
\end{minipage}
\end{minipage}
\end{center}
\vspace{10pt}

A note on how to use the thing. Chapter~\ref{ch:language} is about scepticism, not about UFOs,
and it comes first for the reason given above: the whole book is about signal
integrity, and Chapter~\ref{ch:language} is where we install the instruments. Skip it and the
rest is just stories. Where I do not know something, I have said so. Where I
lack the equipment or the clearance or the money to settle a question, I have
said that too. Those admissions are not gaps in the argument; they are the
argument, and they are the one habit I would most like you to steal from me.

I never called in to Bell's show. I meant to. Ten years of meaning to, and then
one night in April the line went dead for good, and what I had left was a blank
document and about three hundred pages of things I wanted to say to a man who
was very good at letting people say them.

I did, eventually, get on the air --- not with Bell, but with the man who would
later be the one to tell me he was gone. In the autumn before that night in the
lab, the Ufology course at the University of Lethbridge went out live on
\emph{Coast to Coast AM}, and I spent a few minutes talking to George Noory
about what on earth a Canadian university thought it was doing teaching this
subject at all. I have kept the recording.\footnote{The author's live conversation with George Noory on
\emph{Coast to Coast AM}:
\srclink{https://www.dropbox.com/scl/fi/g3kgjxuoejislpwtq6bbm/U-of-L-Ufology-on-Coast-to-Coast-AM.mov?rlkey=r8apwt11bksvihbrxl4en4a99}.}
It is a strange thing to owe a debt to the same broadcast twice: first for the
only three minutes in which I got to answer back, and then, a few months later,
for the worst news I ever heard through headphones.

\par\vspace{7pt}
\begin{center}
\includegraphics[width=0.48\textwidth]{rushweek.jpg}\par
\vspace{4pt}
\begin{minipage}{0.72\textwidth}
{\RaggedRight\centering\scriptsize\color{steel}\itshape The author in 2017, manning the
U~of~L Ufology Group table during Rush Week, with the shirt reading ``Wanna Take A Ride?''
under a banner reading ``Come take a ride on the wild side!'' --- Art Bell's sign-off, hanging
over that table for years before I noticed whose line I had been standing under. Author's own
photograph.\par}
\end{minipage}
\end{center}
% The preface closes on this page, so the closing quotation and the signature are
% pushed to the foot of it with \vfill rather than left sitting in the middle of a
% half-empty page: the quote is the last thing the reader sees before the
% Introduction plate, and no slab of whitespace is left under it.
\vfill
\begin{center}
{\itshape\large ``Come take a ride on the wild side!''}
\end{center}
\vspace{4pt}

\begin{flushright}
\itshape Evan Svendsen
\end{flushright}
\vspace{2pt}

\clearpage

% ===================== INTRODUCTION =====================
\chapterplate{pl_intro.jpg}
\chapter*{Introduction}
\addcontentsline{toc}{chapter}{Introduction}
\markboth{Introduction}{Introduction}

\chapquote{``The absence of evidence is not the evidence of absence.''}{commonly attributed to Carl Sagan}

I am a self-made professor of Ufology. There is no such degree. No university
on this planet will hand you a piece of paper that says you have mastered the
study of unidentified flying objects, and if one ever does, I would advise you
to read the fine print very carefully before you sign the tuition cheque. What
I have instead is a decade of case files, a stack of declassified documents
that cost me more in printer toner than I care to admit, a standing invitation
to lecture at the University of Lethbridge, and a stubborn refusal to accept
``swamp gas'' as a final answer to anything.

This book began as those lectures. Each chapter that follows corresponds
roughly to a week of material, delivered to a room of students who arrived
expecting little green men and left, I hope, with something considerably more
useful: a working method for thinking about extraordinary claims without
either swallowing them whole or spitting them out on reflex.

Which brings me to the single most important structural decision in this
volume, and the one that confuses people most. Chapter~\ref{ch:language} is not about UFOs.
Chapter~\ref{ch:language} is about scepticism. The introduction to Ufology proper does not
arrive until Chapter~\ref{ch:intro}. This ordering is deliberate and it is not negotiable.
Scepticism is a hard dependency. You cannot meaningfully evaluate a single
photograph, a single radar return, or a single witness statement until you
understand how language distorts questions before they are even asked, how
the burden of proof actually works, why correlation is not causation, and how
your own brain will cheerfully lie to you in service of what you already
believe. Attempting Chapter~\ref{ch:ufocsi} without Chapter~\ref{ch:language} is like attempting calculus
without algebra. You will produce answers. They will be wrong, and worse, you
will not know why.

A word on what this book is not. It is not a catalogue of sightings; there
are better catalogues, maintained by people with more filing cabinets than I
have. It is not an argument that aliens are visiting Earth. I do not know
that, you do not know that, and anyone who tells you they know that is
selling something. Nor is it a debunking, because debunking is just believing
with the sign flipped, and I find it every bit as lazy.

What this book is, is an attempt to treat a disreputable subject with
reputable methods. To ask, of each case and each claim: what would count as
evidence here? What test could be run? What would have to be true for this to
be false? Those are scientific questions, and they can be asked about
anything, including things that make you look ridiculous at dinner parties. I
have been looking ridiculous at dinner parties for years. You get used to it.

You will find, scattered through these pages, moments where I tell you
plainly that I do not know something, or that I lack the resources to test it
myself, or that a piece of celebrated evidence collapses the moment you touch
it. Those admissions are not weaknesses in the argument. They are the
argument. A field that cannot say ``I was wrong'' is not a field, it is a
congregation.

So read Chapter~\ref{ch:language} first. Read it twice if you are the sort of person who
already knows what you think about all this. Then we will get to the good
stuff: Area 51, Roswell, Rendlesham, the physics of time travel, the Drake
equation, cattle mutilations, crystal skulls, and the long strange question
of what happens on the day we finally answer the phone.

\bigskip
{\centering
{\small\bfseries This book contains EXACTLY \fmtwordcount{} \booklanguage{} WORDS!}\par}
\medskip
\begin{center}
\begin{minipage}{0.92\textwidth}
Ninety-nine point nine nine per cent of those words were written by artificial intelligence.
The AI did not decide what to say. It did not choose the structure, select the evidence,
or spend a decade building the university course curriculum this textbook is built on.
That was me. The hard work was mine. The words are the machine's.

I am not hiding that, and I am not apologising for it either, because the number itself
makes my point stronger than any argument could: even when you know every word was
written by a machine, you still have to read it. You still have to decide whether each
sentence is true. You still have to sort signal from noise. That is the problem with human
language. It does not get easier just because you know who --- or what --- typed it.
\end{minipage}
\end{center}

